% Generated by scripts/assemble.py; corrections belong in source/corrections.json.
\examhead{41}{同济大学 2026 年数学分析试卷}{44}
% p045-b04; source PDF page 45
\par\Needspace{6\baselineskip}
\blockanchor{p045-b04}
\question{1}(10 分) 判断极限\par
\nopagebreak[4]
% p045-b05; source PDF page 45
\blockanchor{p045-b05}
\displayformula{\lim _ {(x, y) \rightarrow (0, 0)} \frac {\sin (x + y ^ {2}) + \tan (x ^ {2} + y) - x - y}{\arcsin (x ^ {2} + y ^ {2})}}
% p045-b06; source PDF page 45
\blockanchor{p045-b06}
是否存在, 若存在求出极限, 若不存在, 说明理由.\par
% p045-b07; source PDF page 45
\blockanchor{p045-b07}
\question{2}(15 分) 设 $f(x)$ 是 $[a, b]$ 上可导的函数, 证明: 若 $f(x)$ 满足 $|f'(x)| \leq M |f(x)|$ ($M > 0$), 且 $f(a) = 0$, 则 $f(x)$ 恒为常数.\par
% p045-b08; source PDF page 45
\par\Needspace{6\baselineskip}
\blockanchor{p045-b08}
\question{3}(15 分) 设 $f(x) \in C^2[a, b]$，证明梯形求积公式误差满足\par
\nopagebreak[4]
% p045-b09; source PDF page 45
\blockanchor{p045-b09}
\displayformula{\int_ {a} ^ {b} f (x) \mathrm{d} x - \frac {b - a}{2} [ f (a) + f (b) ] = - \frac {(b - a) ^ {3}}{12} f ^ {\prime \prime} (\xi).}
% p045-b10; source PDF page 45
\blockanchor{p045-b10}
其中  $\xi \in [a, b]$ .\par
% p045-b11; source PDF page 45
\blockanchor{p045-b11}
\question{4}(20 分) 设 $T(r)$ 是定义在 $[r_0, +\infty)(r_0 > 0)$ 上的一个单调递增函数, 且 $T(r_0) \geq 1$. 证明:\par
% p045-b12; source PDF page 45
\par\Needspace{6\baselineskip}
\blockanchor{p045-b12}
(1) 若 $T(r) \in C^0[r_0, +\infty)$, 则存在 $\{r_n\} \subset [r_0, +\infty)$ 满足 $r_n \to +\infty$, 使得\par
\nopagebreak[4]
% p045-b13; source PDF page 45
\blockanchor{p045-b13}
\displayformula{T \left(r _ {n} + \frac {1}{T (r _ {n})}\right) <   2 T (r _ {n}).}
% p045-b14; source PDF page 45
\par\Needspace{6\baselineskip}
\blockanchor{p045-b14}
(2) 若 $T(r) \in C^{1}[r_{0}, +\infty)$, 则存在 $\{r_{n}\} \subset [r_{0}, +\infty)$, 满足 $r_{n} \to +\infty$, 使得\par
\nopagebreak[4]
% p045-b15; source PDF page 45
\blockanchor{p045-b15}
\displayformula{T ^ {\prime} (r _ {n}) <   T ^ {1 + \varepsilon} (r _ {n}), \varepsilon > 0.}
% p045-b16; source PDF page 45
\blockanchor{p045-b16}
\question{5}(15 分) 设数域 $\mathbb{F}$ 上的绝对值 $|\cdot|$ 满足 $|x + y| \leq \max \{|x|, |y|\}$, $x, y \in \mathbb{F}$, 而且 $\mathbb{F}$ 关于 $|\cdot|$ 是完备的.\par
% p045-b17; source PDF page 45
\blockanchor{p045-b17}
(1) 试给出 $\mathbb{F}$ 上级数 $\sum_{n=1}^{\infty} a_n$ 收敛的定义.\par
% p045-b18; source PDF page 45
\blockanchor{p045-b18}
(2) 确定 $\sum_{n=1}^{\infty} a_n$ 收敛与 $\lim_{n \to \infty} a_n = 0$ 的关系, 并证明你的结论.\par
% p045-b19; source PDF page 45
\par\Needspace{6\baselineskip}
\blockanchor{p045-b19}
\question{6}(15 分) 设 $f: \mathbb{R}^n \to \mathbb{R}^n$ 是 $C^r$ 映射 $(r \geq 1)$, 若 $f$ 的 Jacobi 行列式在 $x \in \mathbb{R}^n$ 处不为零, 证明: 存在点 $x$ 的开邻域 $U$ 和实数 $\omega > 0$, 使得\par
\nopagebreak[4]
% p045-b20; source PDF page 45
\blockanchor{p045-b20}
\displayformula{\| f (\boldsymbol {x} ^ {\prime}) - f (\boldsymbol {x} ^ {\prime \prime}) \| \geq \omega \| \boldsymbol {x} ^ {\prime} - \boldsymbol {x} ^ {\prime \prime} \|, \forall \boldsymbol {x} ^ {\prime}, \boldsymbol {x} ^ {\prime \prime} \in U.}
% p045-b21; source PDF page 45
\blockanchor{p045-b21}
这里 $\| \cdot \|$ 为向量的2-范数.\par
% p045-b22; source PDF page 45
\par\Needspace{6\baselineskip}
\blockanchor{p045-b22}
\question{7}(15 分) 设开集 $U \subset \mathbb{R}^3$, 函数 $u \in C^2(U)$, 且在 $U$ 中满足 $\frac{\partial^2 u}{\partial x^2} + \frac{\partial^2 u}{\partial y^2} + \frac{\partial^2 u}{\partial z^2} = 0$. 证明: 对任意 $\pmb{x}_0 \in U$, 有\par
\nopagebreak[4]
% p045-b23; source PDF page 45
\blockanchor{p045-b23}
\displayformula{u (\boldsymbol {x} _ {0}) = \frac {1}{| \partial B (\boldsymbol {x} _ {0} , r) |} \iint_ {\partial B (\boldsymbol {x} _ {0}, r)} u \mathrm{d} S.}
% p045-b24; source PDF page 45
\blockanchor{p045-b24}
其中 $B(\pmb{x}_0, r)$ 为包含于 $U$ 中的以 $x_0$ 为球心，$r$ 为半径的任意球，$|\partial B(\pmb{x}_0, r)|$ 为球面面积.\par
% p046-b00; source PDF page 46
\blockanchor{p046-b00}
\question{8}(20 分) 判断下列反常积分的敛散性, 若收敛, 计算积分值; 若发散, 说明理由.\par
% p046-b01; source PDF page 46
\blockanchor{p046-b01}
(1) $\iint_{D} \frac{x^{2} - y^{2}}{(x^{2} + y^{2})^{2}} \mathrm{d}x \mathrm{d}y,$ 其中 $D = \{(x, y) \mid x \geq 1, y \geq 1\}$.\par
% p046-b02; source PDF page 46
\blockanchor{p046-b02}
(2) $\iiint_{D} \frac{\mathrm{d}x\mathrm{d}y\mathrm{d}z}{(x + y + z)^2}$, 其中 $D = \{(x, y, z) \mid x > 0, y > 0, z > 0, x + y + z < 1\}$.\par
% p046-b03; source PDF page 46
\blockanchor{p046-b03}
\question{9}(15 分) 设 $f(x, y) = \int_{0}^{+\infty} \frac{\arctan xt \cdot \arctan yt}{t^2} \, \mathrm{d}t (x > 0, y > 0)$, 计算 $\frac{\partial^2 f}{\partial x \partial y}$.\par
% p046-b04; source PDF page 46
\blockanchor{p046-b04}
\question{10}(10 分) 设  $\boldsymbol{F} = (x - z, x^{3} - yz, -3xy^{2})$ ,  $S: z = \sqrt{4 - x^{2} - y^{2}}$ , n 为 S 上侧的单位法向量, 计算曲面积分  $\iint_{S} \operatorname{rot} \boldsymbol{F} \cdot \boldsymbol{n} \, dS$ .\par
